\documentclass{article} % say
\usepackage{amsmath}
\usepackage{amssymb}
\usepackage{gensymb}
\usepackage{hyperref}
\usepackage{tikz}
\usepackage{tkz-euclide}
\usepackage{listings}
\usetikzlibrary {babel,angles,quotes,positioning}
\usepackage[utf8]{inputenc}
\usepackage[russian]{babel}

\usepackage{mathtools}
\DeclarePairedDelimiter\ceil{\lceil}{\rceil}
\DeclarePairedDelimiter\floor{\lfloor}{\rfloor}
\begin{document}

\textbf{Задача}: На массиве $a$ длины $n$ задана следующая операция: выбирается такой индекс $i$ ($1 \leqslant i \leqslant n$), что $a_i = i$ и массив сдвигается циклически на $i$ шагов влево. Нам дан массив $b$ длины $n$ и число $k$. Нужно проверить, существует ли такой массив, что после применения к нему указанной операции $k$ раз получается массив $b$. Ожидаемая сложность $\mathcal{O}(n)$.

\clearpage

\textbf{Решение}: 

Заметим, что после применения операции число $i$ становится на последнее место в массиве. Таким образом, мы можем ``откатить'' операцию на шаг назад, передвинув число с последнего места на ``его'' место (т.е. сдвинув массив на $a_n$ позиций вправо). Конечно, это возможно, если $1 \leqslant a_n \leqslant n$. Так как всего различных циклических сдвигов массива не больше $n$, то не больше, чем за $n$ ``откатываний'' мы придём в уже известную позицию.

Таким образом, мы делаем ``откатывания'' пока не увидим число, выходящее за пределы размеров массива, либо до зацикливания. В первом случаем необходимо число $k$ сравнить с количеством шагов, которых нам удалось сделать. Во втором случае массив существует для любого $k$.

Вместо того, чтоб действительно делать циклические сдвиги, можно просто хранить суммарный текущий сдвиг.

\lstinputlisting[language=C++]{main.cpp}

\textbf{Источник}: \url{https://codeforces.com/contest/1893/problem/A}

\end{document}


